% Generated draft: requires mathematical and source review.
\documentclass[11pt]{article}
\usepackage[margin=25mm]{geometry}
\usepackage{fontspec}
\setmainfont{Latin Modern Roman}
\usepackage{amsmath,amssymb,mathtools,unicode-math}
\setmathfont{Latin Modern Math}
\usepackage{enumitem,needspace,xurl,hyperref}
\setlength{\parindent}{0pt}
\setlength{\parskip}{5pt}
\setlength{\emergencystretch}{3em}
\newcommand{\N}{\mathbb{N}}
\newcommand{\Z}{\mathbb{Z}}
\newcommand{\Q}{\mathbb{Q}}
\newcommand{\R}{\mathbb{R}}
\newcommand{\C}{\mathbb{C}}
\newcommand{\F}{\mathbb{F}}
\newcommand{\A}{\mathbb{A}}
\newcommand{\PP}{\mathbb{P}}
\newcommand{\E}{\mathbb{E}}
\newcommand{\eps}{\varepsilon}
\newcommand{\dd}{\,\mathrm{d}}
\begin{document}
\section*{UnsolvedMath ID: 800026}
\textbf{Source problem code:} catalogue entry
\subsection{Definitions and mathematical statement}
Let \(X\) be a set, and let \(X^X\) denote the set of all maps from \(X\) to itself, with composition as multiplication. A transformation semigroup on \(X\) is a nonempty subset \(S\subseteq X^X\) closed under composition.

Under the usual interpretation of pointwise topology, give \(X\) the discrete topology and \(X^X\) the product topology. A basic neighbourhood of \(f\in X^X\) is
\[
U(f,F)=\{g\in X^X:g(x)=f(x)\text{ for every }x\in F\},
\qquad F\subseteq X\text{ finite}.
\]
The supplied excerpt does not explicitly specify this discrete topology convention or restrictions on \(X\).

\begin{enumerate}
\item Describe the transformation semigroups \(S\subseteq X^X\) that are closed in the pointwise topology. Under the convention above, the property to be characterized is
\[
\forall f\in X^X,\quad
\left[
\forall F\subseteq X\text{ finite},\ \exists s\in S,\ 
\forall x\in F,\ s(x)=f(x)
\right]
\Longrightarrow f\in S.
\]
The source does not specify the desired form of the description or a class of underlying sets.
\item Describe endomorphism semigroups that are closed in a topology of Zariski type. For a specified mathematical structure \(\mathcal A\), write \(\operatorname{End}(\mathcal A)\) for its semigroup of self-maps preserving its specified structure. The supplied excerpt does not define the topology of Zariski type, identify its ambient space, specify the allowed structures, or clarify whether the objects to be classified are full endomorphism semigroups or subsemigroups of them. Consequently, this second request cannot be stated as a uniquely determined closure condition from the supplied evidence.
\end{enumerate}

The source attributes both requests to G. Mashevitzky. It reports that transformation semigroups closed in the pointwise topology are precisely endomorphism semigroups of mathematical structures, referring to [Mash2]. It also reports that [Mash1] establishes that the topology of Zariski type is weaker than the pointwise topology and gives conditions for their equality, and that [Mash2] resolves several special cases. These are supplied background claims; their hypotheses and the relevant structures and topologies are not specified in the excerpt and are not independently verified here.
\subsection{Short English statement}
Problem 27, attributed to G. Mashevitzky, asks for descriptions of transformation semigroups closed under pointwise limits and of endomorphism semigroups closed in a topology of Zariski type. The first request has the displayed finite interpolation formulation under the discrete topology convention. The second requires definitions absent from the supplied excerpt. The registry labels the problem open; this is a supplied status, not an independently verified assessment.
\subsection{Sources}
\begin{itemize}
\item Local UnsolvedMath registry: permanent ID \(800026\), source problem code \texttt{catalogue entry}, title \emph{The Lyapin's Notebook --- Problem 27}. The supplied complete record labels the problem open and lists Bershadsky, Kublanovsky, and Mashevitzky in its proposer field, with year 2026; the statement itself attributes the problem to G. Mashevitzky.
\item Bershadsky, Kublanovsky, and Mashevitzky, \emph{The Lyapin's Notebook: Unsolved Problems in Semigroup Theory} (2026), Problem 27. Supplied source URL: \url{https://arxiv.org/abs/2604.04763}. The record reports access on 2026-10-16 and extraction by \texttt{arxiv\_latex}; these are supplied provenance details, not a claim of independent source access or verification.
\item [Cameron], [Mash1], [MPP], and [Mash2] are the reference labels supplied for motivation and definitions. The excerpt assigns the comparison of topologies and conditions for equality to [Mash1], and the endomorphism characterization and several special cases to [Mash2]. Full bibliographic entries and links for these four references are not supplied.
\end{itemize}
\end{document}
